% ============================================================
% Chapter 4: Proposed Method: Specification and Design
%   Section: The Fact-Checking Pipeline
%     Subsection: The Second Look
% Included from chapters/04-proposed-method/06-fact-checking-pipeline/fact-checking-pipeline.tex
% Source: backend/src/services/fact_checker/second_look.py;
% docs/decisions/evaluation.md, "The second look".
% ============================================================

\subsection{The Second Look}
\label{subsec:second-look}

The first search for a claim is planned from the claim alone: who and
what it names, what it asserts, and its refutation
(Subsection~\ref{subsec:evidence-retrieval}). When those three queries
find nothing usable the claim ends \texttt{UNVERIFIED}; when they find
something thin, the model answers with low confidence. In both cases the
model's own explanation usually says what was missing --- ``no source
reports the 2024 figure'', ``the sources describe a different
programme''. The second look turns that sentence into a second search.

\begin{enumerate}
    \item \textbf{Follow-up queries.} The model is given the claim, the
    first verdict, its explanation and the titles of the sources it had,
    and asked for at most three queries, in the claim's language, aimed
    at what was missing: the original study or report, an official
    figure, the organisation's own announcement, an independent report
    of the same event. A model that writes none gets the claim itself as
    its query, a sentence the first search never sent whole.
    \item \textbf{Only new pages.} The queries go through the same
    search, with the same concurrency and fusion, and only pages the
    first pass never saw, on domains it did not already have, are kept;
    the article itself never is.
    \item \textbf{One judgement over both.} The first pass's evidence
    and the new pages are ranked, gated, judged by the model and
    recalibrated together, by the same code as the first pass. The
    second verdict replaces the first, which stays on the claim with the
    reason the second look was taken and the queries it sent, so a
    reader can see that the verdict came from a second search.
\end{enumerate}

It fires only when a run enables it, and only on a verdict in doubt:
\texttt{UNVERIFIED}, or a confidence below a threshold (0.5 by default).
It never fires on an outage --- with no search or no model, a second
search meets the same outage --- and never twice. It follows the
iterative retrieval of FIRE \cite{xie2025fire} and the question
generation of AVeriTeC's baseline \cite{schlichtkrull2023averitec},
reduced to one extra round: on a processor without a graphics card each
round costs a search, a handful of page fetches and two model calls.

Whether it should be on is not assumed. It is decided on the tuning
half of the custom set (Subsection~\ref{subsec:custom-split}) by the
tuning loop of Subsection~\ref{subsec:tuning-loop}, and measured on the
test half.
